\documentclass[12pt]{article}
\usepackage{amsmath, amssymb}
\usepackage[margin=1in]{geometry}
\setlength{\parskip}{6pt}
\title{QUBO Formulation: Densest K-Subgraph Problem}
\date{}
\begin{document}
\maketitle

\subsection*{Densest K-Subgraph Problem}

\paragraph{Problem Statement.}
Let $G = (V, E)$ be an undirected weighted graph with vertex set $V = \{0, 1, \dots, N-1\}$ and edge set $E$. Each edge $\{i, j\} \in E$ has a non-negative weight $w_{ij}$. The Densest K-Subgraph Problem asks for a subset of vertices $S \subseteq V$ such that $|S| = K$ and the total weight of the edges induced by $S$ is maximized. Formally, we seek to maximize
\begin{equation}
    \sum_{\{i, j\} \subseteq S} w_{ij},
\end{equation}
subject to the cardinality constraint $|S| = K$.

\paragraph{Decision Variables.}
We introduce a binary variable $x_i \in \{0, 1\}$ for each vertex $i \in V$. The variable $x_i$ indicates whether vertex $i$ is included in the selected subset:
\begin{equation}
    x_i = \begin{cases}
        1 & \text{if } i \in S, \\
        0 & \text{otherwise}.
    \end{cases}
\end{equation}
The selection of exactly $K$ vertices is enforced by the constraint $\sum_{i \in V} x_i = K$.

\paragraph{QUBO Derivation.}
We construct a Quadratic Unconstrained Binary Optimization (QUBO) formulation by converting the maximization problem into a minimization problem and incorporating the cardinality constraint as a quadratic penalty term.

\textbf{Step 1: Objective Function.}
Maximizing the induced edge weight is equivalent to minimizing its negative. The objective function in terms of the decision variables is:
\begin{equation}
    \min_{x \in \{0,1\}^N} \left( -\sum_{\{i, j\} \in E} w_{ij} x_i x_j \right).
\end{equation}
For completeness, we assume $w_{ij} = 0$ for non-edges, allowing the sum to range over all pairs $i < j$.

\textbf{Step 2: Constraint.}
The requirement to select exactly $K$ vertices is expressed as:
\begin{equation}
    \sum_{i \in V} x_i = K.
\end{equation}

\textbf{Step 3: Penalty Term.}
We introduce a penalty term with weight $\lambda > 0$ to enforce the constraint. The penalty is the square of the constraint violation:
\begin{equation}
    P(x) = \lambda \left( \sum_{i \in V} x_i - K \right)^2.
\end{equation}
Expanding the square and using the idempotent property of binary variables ($x_i^2 = x_i$), we obtain:
\begin{align}
    P(x) &= \lambda \left( \sum_{i \in V} x_i^2 + \sum_{i \neq j} x_i x_j - 2K \sum_{i \in V} x_i + K^2 \right) \\
    &= \lambda \left( \sum_{i \in V} x_i + 2 \sum_{i < j} x_i x_j - 2K \sum_{i \in V} x_i + K^2 \right) \\
    &= \lambda \left( (1 - 2K) \sum_{i \in V} x_i + 2 \sum_{i < j} x_i x_j + K^2 \right).
\end{align}
Thus, the penalty contributes linear terms with coefficient $\lambda(1 - 2K)$, quadratic terms with coefficient $2\lambda$, and a constant offset $\lambda K^2$.

\textbf{Step 4: Combined Hamiltonian.}
The total QUBO Hamiltonian $H(x)$ is the sum of the negated objective and the penalty term:
\begin{equation}
    H(x) = \sum_{i \in V} \lambda(1 - 2K) x_i + \sum_{i < j} \left( 2\lambda - w_{ij} \right) x_i x_j + \lambda K^2.
\end{equation}

\textbf{Step 5: QUBO Parameters.}
Reading off the coefficients from $H(x)$, the QUBO matrix entries $Q_{i,j}$ and offset $c$ are given by:
\begin{align}
    Q_{i,i} &= \lambda(1 - 2K), \quad \forall i \in V, \\
    Q_{i,j} &= 2\lambda - w_{ij}, \quad \forall i < j, \\
    c &= \lambda K^2.
\end{align}
The QUBO problem is to find $x^* \in \{0,1\}^N$ that minimizes $x^\top Q x + c$.

\paragraph{Penalty Weight Justification.}
The penalty weight $\lambda$ must be chosen sufficiently large to ensure that any infeasible solution has a strictly higher energy than the optimal feasible solution. Let $W_{\max} = \max_{\{i,j\} \in E} w_{ij}$ denote the maximum edge weight.

Consider an infeasible solution with $K+1$ selected vertices. The penalty incurred is $\lambda$. The maximum possible gain in the objective function by selecting an extra vertex, compared to the best feasible $K$-subgraph, is bounded by the maximum degree of the graph times the maximum edge weight. Specifically, removing any vertex from a set of $K+1$ vertices yields a valid $K$-subgraph, and the weight lost is at most $(N-1)W_{\max}$. Thus, the objective value of the $K+1$-subgraph can exceed the optimal objective value by at most $(N-1)W_{\max}$. To penalize this violation, we require:
\begin{equation}
    \lambda > (N - 1) W_{\max}.
\end{equation}
A symmetric argument applies to solutions with $K-1$ vertices. Therefore, choosing $\lambda > (N-1)W_{\max}$ guarantees that the global minimum of $H(x)$ corresponds to a feasible solution with maximum induced weight.

\paragraph{Correctness Argument.}
Minimizing $H(x)$ over $\{0,1\}^N$ recovers the optimal solution to the Densest K-Subgraph Problem. For any feasible solution $x$ (where $\sum_i x_i = K$), the penalty term vanishes, and $H(x)$ equals the negated objective value. For any infeasible solution $x'$, the penalty term is at least $\lambda$. By the choice of $\lambda$, the energy increase due to the penalty strictly outweighs any potential reduction in the objective term. Consequently, the global minimum of $H(x)$ must occur at a feasible point, and among feasible points, $H(x)$ is minimized exactly when the induced edge weight is maximized.

\paragraph{Illustrative Example.}
Consider a graph with $N=3$ vertices $V=\{0,1,2\}$ and $K=2$. Let the edge weights be $w_{01}=1$, $w_{02}=2$, and $w_{12}=3$. The maximum weight is $W_{\max}=3$. We choose $\lambda = 7$, which satisfies $\lambda > (3-1) \times 3 = 6$.

The QUBO parameters are:
\begin{align}
    Q_{0,0} = Q_{1,1} = Q_{2,2} &= 7(1 - 4) = -21, \\
    Q_{0,1} &= 14 - 1 = 13, \\
    Q_{0,2} &= 14 - 2 = 12, \\
    Q_{1,2} &= 14 - 3 = 11, \\
    c &= 7 \times 2^2 = 28.
\end{align}
The Hamiltonian is:
\begin{equation}
    H(x) = -21(x_0 + x_1 + x_2) + 13 x_0 x_1 + 12 x_0 x_2 + 11 x_1 x_2 + 28.
\end{equation}
Evaluating $H(x)$ for feasible solutions (those with exactly two 1s):
\begin{itemize}
    \item $x = (1,1,0)$: $H = -42 + 13 + 28 = -1$ (induced weight $1$).
    \item $x = (1,0,1)$: $H = -42 + 12 + 28 = -2$ (induced weight $2$).
    \item $x = (0,1,1)$: $H = -42 + 11 + 28 = -3$ (induced weight $3$).
\end{itemize}
The infeasible solution $x=(1,1,1)$ yields $H = -63 + 13 + 12 + 11 + 28 = 1$, which is higher than any feasible energy. The minimum energy is $-3$, achieved at $x^* = (0,1,1)$, corresponding to the optimal vertex subset $\{1,2\}$ with maximum induced weight $3$.

\end{document}